\subsection{ORDER BY দিয়ে ফলাফল সাজানো}

অংশ~\ref{sec:ch06-student-master}-এর একই Student টেবিলে নিচের ফল পাওয়া যায়। ORDER BY-এর নিয়ম Sorting অংশে আছে; এখানে ভিন্ন মানদণ্ডের প্রয়োগ দেখানো হয়েছে।

\noindent\begin{minipage}{\linewidth}
\chSixExample{৭.২২}{Name-এর বর্ণানুক্রমে শিক্ষার্থীদের তথ্য সাজাও।}

\noindent\begin{minipage}[c]{.46\linewidth}
\begin{verbatim}
SELECT Roll, Name, Department, GPA
FROM Student
ORDER BY Name ASC;
\end{verbatim}
\end{minipage}\hfill\begin{minipage}[c]{.04\linewidth}\raggedright{\color{chSixTeal}$\Rightarrow$}\end{minipage}\hfill
\begin{minipage}[c]{.47\linewidth}\raggedright\footnotesize\setlength{\tabcolsep}{2pt}\textbf{ফল}\par\smallskip
\begin{tabular}{|l|l|l|l|}\hline\chSixHeader
Roll & Name & Department & GPA\\\hline
205 & Farhan & Business Studies & 4.80\\\hline
306 & Mahi & Humanities & 4.50\\\hline
150 & Mim & Science & 4.20\\\hline
109 & Nabila & Science & 4.80\\\hline
102 & Rahim & Science & 5.00\\\hline
305 & Rima & Humanities & 5.00\\\hline
213 & Sami & Business Studies & 5.00\\\hline
180 & Tania & Business Studies & 3.90\\
\hline\end{tabular}
\end{minipage}
\end{minipage}\par\smallskip

\noindent\begin{minipage}{\linewidth}
\chSixExample{৭.২৩}{GPA বড় থেকে ছোট ক্রমে সাজাও।}

\noindent\begin{minipage}[c]{.53\linewidth}
\begin{verbatim}
SELECT Roll, Name, GPA
FROM Student
ORDER BY GPA DESC;
\end{verbatim}
\end{minipage}\hfill\begin{minipage}[c]{.04\linewidth}\raggedright{\color{chSixTeal}$\Rightarrow$}\end{minipage}\hfill
\begin{minipage}[c]{.40\linewidth}\raggedright\footnotesize\setlength{\tabcolsep}{2pt}\textbf{ফল}\par\smallskip
\begin{tabular}{|l|l|l|}\hline\chSixHeader
Roll & Name & GPA\\\hline
102 & Rahim & 5.00\\\hline
213 & Sami & 5.00\\\hline
305 & Rima & 5.00\\\hline
109 & Nabila & 4.80\\\hline
205 & Farhan & 4.80\\\hline
306 & Mahi & 4.50\\\hline
150 & Mim & 4.20\\\hline
180 & Tania & 3.90\\
\hline\end{tabular}
\end{minipage}
\end{minipage}\par\smallskip

\noindent\begin{minipage}{\linewidth}
\chSixExample{৭.২৪}{GPA বড় থেকে ছোট সাজাও; GPA সমান হলে Roll ছোট থেকে বড় সাজাও।}

\noindent\begin{minipage}[c]{.53\linewidth}
\begin{verbatim}
SELECT Roll, Name, GPA
FROM Student
ORDER BY GPA DESC, Roll ASC;
\end{verbatim}
\end{minipage}\hfill\begin{minipage}[c]{.04\linewidth}\raggedright{\color{chSixTeal}$\Rightarrow$}\end{minipage}\hfill
\begin{minipage}[c]{.40\linewidth}\raggedright\footnotesize\setlength{\tabcolsep}{2pt}\textbf{ফল}\par\smallskip
\begin{tabular}{|l|l|l|}\hline\chSixHeader
Roll & Name & GPA\\\hline
102 & Rahim & 5.00\\\hline
213 & Sami & 5.00\\\hline
305 & Rima & 5.00\\\hline
109 & Nabila & 4.80\\\hline
205 & Farhan & 4.80\\\hline
306 & Mahi & 4.50\\\hline
150 & Mim & 4.20\\\hline
180 & Tania & 3.90\\
\hline\end{tabular}
\end{minipage}
\end{minipage}\par\smallskip

\begin{boardbox}{ORDER BY-এর গুরুত্বপূর্ণ তথ্য}

\begin{itemize}

\item ASC = Ascending।

\item DESC = Descending।

\item ASC/DESC না লিখলে সাধারণত ঊর্ধ্বক্রম ব্যবহৃত হয়।

\item সাজানোর জন্য একাধিক ফিল্ড থাকলে বামদিকের ফিল্ড আগে বিবেচিত হয়।

\item শুধু \texttt{ORDER BY GPA DESC} দিলে সমান GPA রেকর্ডের
নিজেদের মধ্যকার নির্দিষ্ট ক্রম নিশ্চিত নয়।

\end{itemize}

\end{boardbox}

